\chapter{Analisi}

\section{Analisi dei requisiti}

Si vuole realizzare un database a supporto della gestione di un ricettario online, il quale mette a disposizione la possibilità di acquisto. 
La base di dati dovrà quindi immagazzinare informazioni relative agli utenti, alle ricette da loro pubblicate, nonché agli ordini effettuati 
sul sito. Gli amministratori potranno dunque consultare tutte le ricette inserite, aggiungere nuovi ingredienti disponibili e lanciare nuove 
promozioni.


\textbf{\\Intervista}


Si vuole tenere traccia degli utenti registrati al sito memorizzandone nome, cognome, codice fiscale, email, password e opzionalmente un indirizzo di spedizione.
In seguito alla registrazione l’utente può visualizzare e/o creare delle ricette.
\\Chi crea una ricetta vi deve associare un nome, una descrizione della preparazione, un tempo stimato di preparazione e i necessari ingredienti con le relative quantità (espresse in grammi).
\\Gli ingredienti disponibili al momento della creazione della ricetta sono quelli presenti sul sito, ognuno associato ad un codice identificativo, un nome e un relativo prezzo al chilo.
\\Inoltre nel creare la ricetta, l’utente deve associarla a una o più categorie, anche queste visualizzabili nel sito: ogni categoria è costituita da un codice univoco, una descrizione e un nome (vegano, vegetariano, italiano, asiatico, ...).
\\Gli utenti possono per di più lasciare delle recensioni alle ricette sul sito, indicando un voto in una scala da 1 a 10 e un eventuale consiglio di modifica (un utente può scrivere un solo giudizio per una determinata ricetta).
\\In aggiunta, il sistema mette a disposizione la possibilità di acquisto: un utente che visualizza una ricetta può acquistare tutto il necessario per prepararla a casa (eventualmente un utente può acquistare più di una ricetta nello stesso ordine).
\\Il costo di una ricetta è calcolato quindi come somma dei costi degli ingredienti valutandone il peso.
Per semplicità non viene gestita la disponibilità degli ingredienti in magazzino.
\\Infine il sito può lanciare delle promozioni periodiche che offrono degli sconti: in particolare la percentuale di sconto riferita ad una specifica promo dipende dalla categoria di ricetta.
Nel caso una ricetta preveda più sconti al momento dell’acquisto (più categorie scontate), verrà applicato quello maggiore.


\textbf{\\Revisione}

A seguito dell’analisi dell’intervista e del confronto del dominio con applicazioni e siti web realmente esistenti, si è deciso di apportare alcuni cambiamenti nella modellazione e gestione degli sconti, in modo tale da renderle più complesse e realistiche. 
\\\\Una prima modifica proibisce all’utente di pubblicare due ricette con lo stesso nome, in modo tale da prevenire duplicazioni di ricette sul sito.
\\\\Una seconda modifica è riassunta dalla seguente descrizione: “la percentuale di sconto riferita ad una specifica promo dipende dalla categoria di ricetta e dal numero di ingredienti di quest’ultima”.
Nello specifico, uno sconto si riferisce ad un certo intervallo di numero di ingredienti, e una ricetta per essere scontata deve avere un numero di ingredienti appartenente a questo range, oltre che appartenere alla specifica categoria interessata.
\\\\Infine il dominio impone il seguente vincolo: i range di ingredienti di eventuali sconti riferiti ad una stessa categoria e stessa promozione non devono essere sovrapposti.

A seguito delle revisioni si riformula l’intervista nella seguente:

\begin{table}[H]
\centering
\setlength{\extrarowheight}{2pt}
\begin{tabularx}{\textwidth}{|c|X|}
\hline
\multicolumn{2}{|c|}{\textit{\textbf{Definizione delle specifiche in linguaggio naturale}}} \\
\hline
1 & Si vuole tenere traccia degli utenti registrati al sito memorizzando nome, cognome,\\
2 & codice fiscale, email, password e opzionalmente un indirizzo di spedizione. \\
3 & In seguito alla registrazione l'utente può visualizzare e/o creare delle ricette. \\
4 & Chi crea una ricetta vi deve associare un nome, una descrizione della preparazione, \\
5 & un tempo stimato di preparazione e i necessari ingredienti con le relative quantità \\
6 & (espresse in grammi). \\
7 & Gli ingredienti disponibili al momento della creazione della ricetta sono quelli \\
8 & presenti sul sito, ognuno associato ad un codice identificativo, un nome e un relativo \\
9 & prezzo al chilo. \\
10 & Inoltre nel creare la ricetta, l'utente deve associarla a una o più categorie, anche \\
11 & queste visualizzabili nel sito: ogni categoria è costituita da un codice univoco, una \\
12 & descrizione e un nome (vegano, vegetariano, italiano, asiatico, ...). \\
13 & Il sistema proibisce all'utente di pubblicare più ricette con lo stesso nome. \\
14 & Gli utenti possono per di più lasciare delle recensioni alle ricette sul sito, indicando \\
15 & un voto in una scala da 1 a 10 e un eventuale consiglio di modifica (un utente può \\
16 & scrivere un solo giudizio per una determinata ricetta). \\
17 & In aggiunta, il sistema mette a disposizione la possibilità di acquisto: un utente \\
18 & che visualizza una ricetta può acquistare tutto il necessario per prepararla a casa \\
19 & (eventualmente un utente può acquistare più di una ricetta nello stesso ordine). \\
20 & Il costo di una ricetta è calcolato quindi come somma dei costi degli ingredienti \\
21 & valutandone il peso. \\
22 & Per semplicità non viene gestita la disponibilità degli ingredienti in magazzino. \\
23 & Infine il sito può lanciare delle promozioni periodiche che offrono degli sconti: \\
24 & in particolare la percentuale di sconto riferita ad una specifica promo dipende \\
25 & dalla categoria di ricetta e dal numero di ingredienti di quest'ultima: un singolo \\
26 & sconto copre un certo range di numero di ingredienti, e i range di ingredienti di \\
27 & eventuali sconti riferiti ad una stessa categoria e stessa promozione non devono \\
28 & essere sovrapposti. \\
29 & Nel caso una ricetta preveda più sconti al momento dell'acquisto (più categorie \\
30 & scontate e/o più sconti per una singola categoria), verrà applicato quello maggiore. \\
\hline
\end{tabularx}
\caption*{\textbf{Tabella 1.1.1} - \textit{Specifiche di progetto in linguaggio naturale.}}
\end{table}

Si procede con l'individuare i termini principali che riassumono al meglio i concetti del dominio.

\begin{table}[H]
\centering
\renewcommand{\arraystretch}{1.3} % Mantiene un po' di respiro tra le righe
\setlength{\tabcolsep}{8pt}

\begin{tabularx}{\textwidth}{|l|X|}
\hline
\textbf{Termine} & \textbf{Breve descrizione} \\
\hline
Utente & Colui che naviga sul sito e che può visualizzare, pubblicare, recensire e acquistare ricette. \\
Ricetta & Elemento acquistabile che rappresenta un piatto, pensato per essere preparato e consumato. \\
Ingrediente & Oggetto necessario per la preparazione di una ricetta. \\
Categoria & Gruppo di appartenenza di una ricetta secondo un determinato criterio. \\
Recensione & Feedback virtuale riguardante una ricetta e pubblicato da un utente. \\
Ordine & Richiesta d’acquisto che raggruppa una o più ricette selezionate. \\
Promozione & Evento commerciale periodico vantaggioso per gli utenti che desiderano acquistare ricette. \\
Sconto & Ribasso del prezzo di una ricetta secondo certi requisiti. \\
\hline
\end{tabularx}
\caption*{\textbf{Tabella 1.1.2} - \textit{Glossario dei termini.}}
\end{table}


Si tiene conto delle seguenti correzioni di ambiguità:

\begin{table}[H]
\centering
\renewcommand{\arraystretch}{1.3} % Mantiene lo spazio verticale comodo per il testo
\setlength{\tabcolsep}{8pt}      % Distanzia il testo dai bordi laterali

\begin{tabularx}{\textwidth}{|c|l|X|}
\hline
\textbf{Riga} & \textbf{Termine} & \textbf{Nuovo termine} \\
\hline
16 & Giudizio & Recensione \\
18 & …tutto il necessario… & Ingredienti \\
24 & Promo & Promozione \\
\hline
\end{tabularx}
\caption*{\textbf{Tabella 1.1.3} - \textit{Rilevamento delle ambiguità e correzioni proposte.}}
\end{table}


\begin{table}[H]
\centering
\setlength{\extrarowheight}{2pt}
\begin{tabularx}{\textwidth}{|c|X|}
\hline
\multicolumn{2}{|c|}{\textit{\textbf{Secifiche ristrutturate}}} \\
\hline
1 & Degli \textbf{utenti} registrati al sito si memorizzano nome, cognome, codice fiscale, email, \\
2 & password e opzionalmente un indirizzo di spedizione. \\
3 & L’utente può visualizzare e/o creare delle ricette. \\
4 & Chi crea una \textbf{ricetta} vi deve associare un nome, una descrizione della preparazione, \\
5 & un tempo stimato di preparazione e ogni \textbf{ingrediente} necessario con le relative \\
6 & quantità (espresse in grammi). \\
7 & Gli ingredienti disponibili al momento della creazione della ricetta sono quelli \\
8 & presenti sul sito, ognuno associato ad un codice identificativo, un nome e un relativo \\
9 & prezzo al chilo. \\
10 & Inoltre nel creare la ricetta, l’utente deve associarla a una o più categorie, anche \\
11 & queste visualizzabili nel sito: ogni \textbf{categoria} è costituita da un codice univoco, \\
12 & una descrizione e un nome (vegano, vegetariano, italiano, asiatico, ...). \\
13 & Il sistema proibisce all’utente di pubblicare più ricette con lo stesso nome. \\
14 & Per una determinata ricetta presente sul sito, un utente può pubblicare al massimo \\
15 & una recensione, indicando un voto in una scala da 1 a 10 e un eventuale consiglio \\
16 & di modifica. \\
17 & Un utente che visualizza una ricetta può inoltre acquistare l’insieme di tutti gli \\
18 & ingredienti per prepararla a casa (eventualmente un utente può acquistare gli \\
19 & ingredienti di più ricette nello stesso ordine). \\
20 & Il costo di una ricetta è calcolato quindi come somma dei costi degli ingredienti \\
21 & valutandone il peso. \\
22 & Per semplicità non viene gestita la disponibilità degli ingredienti in magazzino. \\
23 & Infine il sito può lanciare delle \textbf{promozioni} periodiche che offrono degli \textbf{sconti}:\\
24 & in particolare la percentuale di sconto riferita ad una specifica promo dipende \\
25 & dalla categoria di ricetta e dal numero di ingredienti di quest’ultima: un singolo \\
26 & sconto copre un certo range di numero di ingredienti, e i range di ingredienti di \\
27 & eventuali sconti riferiti ad una stessa categoria e stessa promozione non devono \\
28 & essere sovrapposti. \\
29 & Nel caso una ricetta preveda più sconti al momento dell’acquisto (più categorie \\
30 & scontate e/o più sconti per una singola categoria), verrà applicato quello maggiore. \\
\hline
\end{tabularx}
\caption*{\textbf{Tabella 1.1.4} - \textit{Specifiche di progetto ristrutturate ed estrazione dei concetti principali.}}
\end{table}


\section{Analisi delle viste}

A seguito delle specifiche ristrutturate, si divide il DataBase in due tipologie di utenti diverse:

\begin{itemize}
    \item \textbf{Utente standard}
    \item \textbf{Amministratore}
\end{itemize}

\subsection{Utenti standard}

I requisiti relativi agli utenti standard sono stati già descritti nel paragrafo precedente e non necessitano di alcuna aggiunta.

\textbf{Elenco delle principali azioni per gli utenti standard:}\\

\begin{enumerate}[label=\textbf{U\arabic* - }, leftmargin=1.2cm, itemsep=0.8em]

    \item REGISTRAZIONE UTENTE\\
    L’utente deve essere in grado di registrarsi correttamente al sito.

    \item VISUALIZZAZIONE DELLE RICETTE\\
    L’utente deve poter essere in grado di visualizzare l’elenco di tutte le ricette pubblicate sul sito.

    \item VISUALIZZAZIONE DI RICETTE PER FILTRO\\
    L’utente deve poter visualizzare delle ricette secondo dei filtri di ricerca adeguati:
    \begin{enumerate}[label=\textbf{U3.\arabic* - }, leftmargin=1.2cm, itemsep=0.2em, topsep=0.3em]
        \item Filtro per nome della ricetta
        \item Filtro per ingrediente
        \item Filtro per tempi di preparazione inferiori a un dato limite
        \item Filtro per prezzo
        \item Filtro per utente
        \item Filtro per categoria/e
    \end{enumerate}

    \item VISUALIZZAZIONE INGREDIENTI\\
    L’utente può visualizzare gli ingredienti disponibili sul sito che possono essere usati per comporre le ricette.

    \item PUBBLICAZIONE DI UNA RECENSIONE\\
    L’utente può valutare una data ricetta pubblicando una recensione contenente voto da 1 a 10 e opzionalmente un commento/suggerimento di modifica.

    \item PUBBLICAZIONE DI UNA RICETTA\\
    L’utente deve poter essere in grado di pubblicare una ricetta sul sito, selezionando gli ingredienti necessari per realizzarla, la/le categoria/e a cui appartiene e i vari attributi.

    \item VISUALIZZAZIONE VANTAGGI\\
    L’utente può visualizzare le categorie di ricette per cui ci sono sconti attivi.

    \item VISUALIZZAZIONE ECCELLENZE\\
    L’utente può visualizzare:
    \begin{enumerate}[label=\textbf{U8.\arabic* - }, leftmargin=1.2cm, itemsep=0.2em, topsep=0.3em]
        \item Le migliori ricette in base alle recensioni
        \item I migliori utenti in base alla media delle recensioni delle ricette da loro pubblicate
        \item Le ricette maggiormente ordinate
        \item Le categorie di ricette maggiormente ordinate
    \end{enumerate}

    \item EFFETTUAZIONE DI UN ORDINE\\
    L’utente può richiedere un ordine, il quale contiene il set di ingredienti di una o più ricette. All’ordine vengono applicati eventuali sconti (ad una ricetta che ha più sconti possibili verrà applicato quello massimo).

    \item VISUALIZZAZIONE STORICO ORDINI\\
    L’utente può visualizzare tutti i suoi ordini effettuati in passato negli ultimi $n$ mesi, in base alle sue preferenze.

    \item VISUALIZZAZIONE SPECIFICA RICETTA\\
    L’utente può visualizzare tutti i dettagli di una specifica ricetta, come il prezzo, l’utente che l’ha pubblicata, la media delle recensioni, l’elenco degli ingredienti e le categorie di appartenenza.

\end{enumerate}

\subsection{Amministratore}

L’amministratore è colui che si occupa di gestire il sito e i relativi dati, preoccupandosi di registrare i nuovi utenti con le relative informazioni.\\
E’ responsabile della gestione del catalogo di ingredienti disponibili, deve quindi poter inserire nuovi ingredienti determinandone il prezzo al chilo.\\
L’amministratore deve inoltre poter inserire nuove categorie di ricetta in base allo sviluppo culinario nel mondo.\\
Deve inoltre tenere sotto controllo l’equilibrio economico del sito, aggiungendo opportunamente nuove promozioni e sconti associati.\\
Infine l’amministratore deve mantenere il sito pulito e coerente: egli ha infatti il potere di rimuovere utenti che hanno pubblicato ricette con dati inadeguati.

\textbf{Elenco delle principali azioni per l'amministratore:}\\

\begin{enumerate}[label=\textbf{A\arabic* - }, leftmargin=1.2cm, itemsep=0.8em]

    \item INSERIMENTO INGREDIENTE\\
    L’amministratore deve poter aggiungere un ingrediente alla lista disponibile, specificando il prezzo al chilo.

    \item AGGIORNAMENTO INGREDIENTE\\
    L’amministratore può modificare il prezzo (al chilo) per un ingrediente disponibile sul sito.

    \item INSERIMENTO CATEGORIA\\
    L’amministratore può aggiungere nuove categorie di ricette disponibili sul sito.

    \item LANCIO DI UNA PROMOZIONE\\
    L’amministratore può inserire nuove promozioni periodiche che annunciano l’inizio di un periodo vantaggioso.

    \item PUBBLICAZIONE SCONTI\\
    L’amministratore può inserire nuovi sconti in base alla promozione di riferimento e alla categoria di ricetta.

    \item RIMOZIONE UTENTI\\
    L’amministratore può rimuovere manualmente utenti che hanno pubblicato diverse ricette con categorie inadeguate. \\
    Il sistema pone particolare attenzione alle ricette vegane e l’amministratore è in grado di individuare in modo automatico ricette classificate come tali ma che non lo sono (ad esempio una ricetta classificata come vegana ma che contiene carne o pesce). In questo particolare caso delle ricette vegane, l’utente verrà rimosso dal sistema alla pubblicazione di almeno 3 ricette inadeguate (e tutte le sue ricette con lui).

\end{enumerate}